% Pursuing Stacks (cotes 134-2 to 134-8): archive cote 134-8, site page 9 (PDF page 10), typescript page 563.
% typescript. Draft compared: ps7.tex:405-485.
% AI-transcribed by Claude Opus 5.5 (claude-opus-5-5, adaptive thinking, effort high), Message Batches API, batch ps4 (msgbatch_012Aoc1NuskATPsyBeKSYK7J), 2026-11-12; reviewed against both supplied scan halves; see CHECK_DECISIONS.tsv.
\scanpage{134-8/9}
\mspage{563}

\margin{\textcircled{$\varphi_A$}}
\diagram{
% Hand-drawn arrows are marked blue; typed objects remain plain.
\begin{tikzcd}[column sep=large]
A^{\hand{\wedge}} \arrow[color=blue,r, "\hand{\phi}_A"] \arrow[color=blue,d, "Wh_A"'] \& (Hot) \arrow[color=blue,d, "\underline{Wh}"] \\
A^{\hand{\wedge}}_{ab} \arrow[color=blue,r, "LH.(A{,}-)"'] \& (Hotab) = D.(Ab)
\end{tikzcd}
}
(7) \hfill ,

or equivalently,  of the corresponding diagram  where the categories
$A^{\hand{\wedge}}$, $A^{\hand{\wedge}}_{ab}$  are replaced by the relevant localizqtions :

\diagram{
\begin{tikzcd}
Hot_A \arrow[color=blue,r] \arrow[color=blue,d] \& Hot \arrow[color=blue,d] \\
Hotab_A \arrow[color=blue,r] \& H\overset{\unsure{o|u}}{\underset{\unsure{o}}{o}}tab
\end{tikzcd}
}
(7') \hfill .

The left vertical arrow in (7)  is of course the trivial abelianization
functor \hand{in $A^{\wedge}$} :
\[ Wh_A :\ A^{\hand{\wedge}} \hand{\longrightarrow} A^{\hand{\wedge}}_{ab} \quad , \qquad X \hand{\longmapsto} \underline{Z}^{(X)} = (\ a \hand{\mapsto} \underline{Z}^{(X(a))}\ ) \quad . \]
\over{g}{G}oing back to the definitions,  the commutativity of (7) up to isomor-%
phism,  means that for X in $A^{\hand{\wedge}}$,  there is a canonical isomorphism
\[ (8) \qquad LH.(A_{/X}\ ,\ \underline{Z}\ ) \hand{\simeq} LH.(A,\ \underline{Z}^{(X)}) \quad , \]
functorial with respect to X.   To define (8),  let
\[ \deleted{f_X\ \text{or }}f :\ A' \overset{def}{=} A_{/X} \hand{\longrightarrow} A \]
be the canonical functor,  then we get (using (*) of previous page,
with the roles of A and A' reversed)
\[ LH.(A',\ \underline{Z}_{A'}) \hand{\simeq} LH.(A,\ Lf^{ab}_{!}(\underline{Z}_{A'})) \quad , \]
and the relation (8) follows from the more precise relation
\[ (9) \qquad Lf^{ab}_{!}(\underline{Z}_{A'}) \xrightarrow{\hand{\sim}} \underline{Z}^{(X)} \quad . \]
To get \xout{the} (9),  we remark that \inserted[typed]{we have the tautological relation} \xout{the right hand side is just}
\[ (9') \qquad f^{ab}_{!}(\underline{Z}_{A'}) \hand{\simeq} \underline{Z}^{(X)} \quad , \]
hence the \xout{arro} map (9).  To prove it is an isomorphism amounts to pro-%
ving
\[ \xout{(10)} \qquad L_i f^{ab}_{!}(\underline{Z}_{A'}) = 0 \quad \text{for} \ \ i \hand{>} 0 \ , \]
but we have indeed
\[ (10) \qquad L_i f^{ab}_{!} \ = 0 \quad \text{for} \ \ i \hand{>} 0 \ , \ \ \text{i.e.} \ \ f^{ab}_{!} \ \ \text{is \uline{exact}} \quad , \]
not only right exact,  a rather special feature\hand{,} \inserted{valid} for a localization

\margin{9}

% Visual collation: both supplied image halves inspected; see CHECK_DECISIONS_134_7_8.tsv.
